\documentclass[11pt]{article}

\usepackage[margin=1in]{geometry}
\usepackage{amsmath,amssymb,amsthm,mathtools}
\usepackage{enumitem}
\usepackage{microtype}
\usepackage{xcolor}
\usepackage{hyperref}

\hypersetup{
  colorlinks=true,
  linkcolor=blue!60!black,
  citecolor=blue!60!black,
  urlcolor=blue!60!black
}

\newtheorem{theorem}{Theorem}[section]
\newtheorem{lemma}[theorem]{Lemma}
\newtheorem{proposition}[theorem]{Proposition}
\newtheorem{corollary}[theorem]{Corollary}
\newtheorem{remark}[theorem]{Remark}

\newcommand{\one}{\mathbf 1}
\newcommand{\Tr}{\operatorname{Tr}}
\newcommand{\ip}[2]{\left\langle #1,#2\right\rangle}
\newcommand{\norm}[1]{\left\lVert #1\right\rVert}
\newcommand{\pos}[1]{\left(#1\right)_+}
\newcommand{\cL}{\mathcal L}
\newcommand{\cS}{\mathcal S}

\title{Region II of the Goodman--Style Odd-Cycle Bound:\\
       a computer-free scalar proof}
\author{Working note for T.K., H.Y., and collaborators}
\date{July 17, 2026}

\begin{document}
\maketitle

\begin{abstract}
Let $W$ be a graphon of edge density $p$, put $q=1-p$, and let $m\ge15$ be
odd.  We prove, throughout the remaining density range
$\tfrac12<p<\tfrac23$, that
\[
        t(C_m,W)\ge p^m-pq^{m-1}.
\]
The operator and eigenfunction argument from the previous Region-II note
reduces the problem to a three-variable Huber payment inequality.  The old
proof of that inequality divided the parameter space into three externally
chosen zones and used two exhaustive interval-box certificates.  Here the
scalar inequality is proved analytically.  At theorem level there is only one
split, namely the intrinsic active-set split
$2\rho\xi\le1$ or $2\rho\xi>1$ in the dual Huber problem.  The first branch is
settled by a quadratic payment estimate and the second by a linear payment
estimate.  All remaining estimates use finite geometric sums, Bernoulli's
inequality, a symmetric second-difference bound, and one-variable calculus.
No numerical sampling, interval arithmetic, or generated certificate is used.

This note replaces the computation-assisted scalar part for odd $m\ge15$.
The finitely many cycles $C_5,\ldots,C_{13}$ remain logically separate and
are not claimed here to have acquired new computer-free proofs.
\end{abstract}

\tableofcontents

\section{Result and proof architecture}\label{sec:intro}

For a graphon $W:[0,1]^2\to[0,1]$, write
\[
        p=t(K_2,W),\qquad q=1-p.
\]
The conjectured Goodman--style inequality for an odd cycle is
\begin{equation}\label{eq:goal}
        t(C_m,W)\ge p^m-pq^{m-1}.
\end{equation}
For $p\le\tfrac12$ the right-hand side is nonpositive.  The present note is
concerned with the nontrivial lower-density band
\[
        \frac12<p<\frac23,
        \qquad \frac13<q<\frac12.
\]

\begin{theorem}[Computer-free Region II for long odd cycles]\label{thm:main}
Let $W$ be a graphon with edge density $\tfrac12<p<\tfrac23$.  For every odd
$m\ge15$,
\[
        t(C_m,W)\ge p^m-p(1-p)^{m-1}.
\]
\end{theorem}

The proof has three layers.  Section~\ref{sec:reduction} records the exact
one-frontier and Huber reduction.  Section~\ref{sec:chart} introduces the
dimensionless chart
\[
        a=\frac q\alpha,\qquad r=\frac Lq,\qquad
        \ell=\frac L\alpha=ar,\qquad N=m-2.
\]
In this chart the Huber dual has two canonical witnesses.  They lead to the
following two analytic statements.

\begin{proposition}[Quadratic certificate]\label{prop:quadratic}
On the branch $2\rho\xi\le1$, the scalar defect satisfies
\[
        R_m\le C_m\psi(\xi,\rho).
\]
\end{proposition}

\begin{proposition}[Linear certificate]\label{prop:linear}
On the branch $2\rho\xi>1$, the scalar defect satisfies
\[
        R_m\le C_m\psi(\xi,\rho).
\]
\end{proposition}

The proofs of these propositions are given in
Sections~\ref{sec:quadratic} and~\ref{sec:linear}.  Some elementary estimates
inside those proofs are separated according to whether a geometric ratio is
close to $1$.  Those are estimates within a lemma, not different graphon
regimes: the only case distinction propagated through the proof is the
active-set distinction $2\rho\xi\le1$ versus $2\rho\xi>1$.

\begin{remark}[What has been removed]\label{rem:removed}
The proof below does not use the former cuts $e=1/60$ and $\xi=1$, the
frontier-gap parameter $\kappa$, an $m$-maximisation by interval bounds, the
secant gate near the Tur\'an corner, or either of the Zone-B and Zone-C box
certifiers.  In particular, the scalar proof can be checked line by line
without running a script.
\end{remark}

\begin{remark}[Why the split at \(p=2/3\) is structural]
The division between Regions I and II should not itself be removed.  In
complement variables it is the threshold \(q=1/3\): the Hilbert--Schmidt
budget forces at most one compression eigenvalue above \(q\) exactly when
\(2q^2\ge q(1-q)\), i.e.\ \(q\ge1/3\).  Below that threshold several frontier
modes are compatible with the same budget.  What the present argument removes
is the downstream zone decomposition inside the one-frontier region.
\end{remark}

\section{The exact one-frontier reduction}\label{sec:reduction}

This section condenses the operator and Huber reduction from
\texttt{paper\_region2\_v2.tex}.  It is included to fix the interface used by
the new scalar proof.

Let $U=1-W$, let $T_U$ be its integral operator, and decompose
$L^2([0,1])=\mathbb R\one\oplus\one^\perp$.  In block form,
\begin{equation}\label{eq:block}
        T_U=
        \begin{pmatrix}q&g^*\\ g&A\end{pmatrix},
        \qquad
        g=T_U\one-q\one,
        \qquad
        A=P_{\one^\perp}T_UP_{\one^\perp}.
\end{equation}
Since $0\le U\le1$,
\begin{equation}\label{eq:HS}
        \Tr(A^2)+2\norm g_2^2\le pq.
\end{equation}
The one-sided determinant expansion gives, for odd $m$,
\begin{equation}\label{eq:shift}
 t(C_m,W)=p^m-\Tr(A^m)+m[z^m]\cL_W(z),
 \qquad
 \cL_W(z)=-\log\!\left(
 1-\frac{z^2\ip g{(I+zA)^{-1}g}}{1-pz}
 \right).
\end{equation}
Indeed, after conjugating by the sign flip on $\one^\perp$, the relevant
operator is
$\left(\begin{smallmatrix}p&g^*\\g&-A\end{smallmatrix}\right)$; Schur's
determinant formula and
$-\log\det(I-zT)=\sum_{j\ge1}z^j\Tr(T^j)/j$ give
\eqref{eq:shift}.  Step-graphon approximation extends the identity from
finite rank to every graphon.

The compression satisfies $\norm A\le\tfrac12$.  If
$\lambda_{\max}(A)\le q$, then
$\Tr(A^m)\le q^{m-2}\Tr(A^2)\le pq^{m-1}$, while the last term in
\eqref{eq:shift} is nonnegative.  Thus only a frontier eigenvalue
$\alpha>q$ needs attention.  It is unique: two such eigenvalues would
contribute more than $2q^2>q(1-q)$ to \eqref{eq:HS}.  Every other eigenvalue
lies in $[-L,L]$, where
\begin{equation}\label{eq:L}
        L=\sqrt{pq-\alpha^2},
        \qquad 0<L<q<\alpha<\frac12.
\end{equation}

For completeness, the eigenfunction shape argument also gives
\begin{equation}\label{eq:forced}
 \norm g_2^2\ge
 \frac{\alpha(\alpha-q)^2}{2(1-2\alpha)},
 \qquad
 \alpha^2+q\alpha-q\le0.
\end{equation}
To see this, choose a unit frontier eigenfunction $\phi\perp\one$ and put
$a_\phi=\int|\phi|$.  Positivity and boundedness of $U$ imply
$a_\phi^2\ge2\alpha$.  Writing
$|\phi|=a_\phi\one+h$, one obtains
\[
 (\alpha-q)a_\phi^2
 \le2a_\phi\norm g_2\sqrt{1-a_\phi^2},
\]
which proves the first inequality in \eqref{eq:forced}.  Combining it with
$\alpha^2+2\norm g_2^2\le pq$ yields the second.  Hence
\begin{equation}\label{eq:ceiling}
        q<\alpha\le r(q):=\frac{\sqrt{q^2+4q}-q}{2}.
\end{equation}

For odd $m$, define
\begin{align}
 k_m(\lambda)&=\frac{p^{m-1}-\lambda^{m-1}}{p+\lambda},\label{eq:kdef}\\
 A_m&=2L^{m-2}+m k_m(\alpha),\qquad
 B_m=2L^{m-2}+m k_m(L),\label{eq:AB}\\
 R_m&=\alpha^m+L^m-pq^{m-1}.\label{eq:R}
\end{align}
The function $k_m$ is positive and strictly decreasing on $[0,p)$, so
\begin{equation}\label{eq:A<B}
        0<A_m<B_m.
\end{equation}
Writing $g=\gamma\phi+g_s$, the trace term and the linear term of the logarithm
in \eqref{eq:shift} give the master estimate
\begin{equation}\label{eq:master}
 t(C_m,W)-\bigl(p^m-pq^{m-1}\bigr)
 \ge -R_m+A_m\gamma^2+B_m\norm{g_s}_2^2.
\end{equation}

It remains to eliminate the shape of $\phi$.  Put
\[
 z_\phi=\left(\int|\phi|\right)^2,\qquad
 b_\phi=\ip{|\phi|}{\phi}\ge0,\qquad
 k_\phi=|\phi|-\sqrt{z_\phi}\one-b_\phi\phi,\qquad
 K=\norm{k_\phi}_2^2.
\]
Then $z_\phi+b_\phi^2+K=1$ and $z_\phi\ge2\alpha$.  Direct use of
$0\le U\le1$ gives
\begin{equation}\label{eq:direct-safe}
 \gamma\le
 \frac{z_\phi-2\alpha-2\alpha b_\phi}{2\sqrt{z_\phi}},
 \qquad
 b_\phi\gamma+\ip{g_s}{k_\phi}
 \ge
 \frac{(\alpha-q)z_\phi+(\alpha-L)K}{2\sqrt{z_\phi}}.
\end{equation}
Minimising the quadratic payment in \eqref{eq:master} subject to
\eqref{eq:direct-safe} is exact and leads to the following quantities.  Set
\begin{equation}\label{eq:dfedef}
        d=\alpha-q,\qquad f=\alpha-L,\qquad e=1-2\alpha,
\end{equation}
and
\begin{align}
 C_m&=\frac{B_m f\sqrt{2\alpha}\,e^2}{4\alpha^2},
 &
 \xi&=\frac{4\alpha^2d}{e^2},
 &
 \rho&=\frac{A_m}{B_m}\frac{\sqrt\alpha}{2\sqrt2 f},
 \label{eq:Cxirho}\\
 \psi(\xi,\rho)
 &=\min_{0\le v\le1}
   \left\{\rho v^2+\pos{\xi-v+v^2}\right\}.\label{eq:psi}
\end{align}
The exact shape elimination is
\begin{equation}\label{eq:huber-elim}
        A_m\gamma^2+B_m\norm{g_s}_2^2
        \ge C_m\psi(\xi,\rho).
\end{equation}
Here is a short derivation.  In the nondegenerate case $K>0$, the second
inequality in \eqref{eq:direct-safe} and Cauchy--Schwarz give
\[
 \norm{g_s}_2^2\ge
 \frac1K\pos{
 \frac{dz_\phi+fK}{2\sqrt{z_\phi}}-b_\phi\gamma}^2.
\]
When the first bound in \eqref{eq:direct-safe} is positive, a minimiser may be
taken with $\gamma\ge0$, and it obeys
\[
 b_\phi^2+2\alpha b_\phi+2\gamma\sqrt{2\alpha}\le e.
\]
Thus, with $v=2\gamma\sqrt{2\alpha}/e\in[0,1]$,
\[
 b_\phi\gamma\le
 \frac{e^2v(1-v)}{4\alpha\sqrt{2\alpha}}.
\]
Minimising first over $K$ by AM--GM gives exactly the right side of
\eqref{eq:huber-elim}.  If the first bound in \eqref{eq:direct-safe} is
nonpositive, the safe-channel payment is at least $B_mdf$, which is stronger.
The case $K=0$ follows directly from the displayed budget (or by its limiting
form): it forces $\xi\le v-v^2$, so the positive-part term vanishes.  This
proves \eqref{eq:huber-elim} in all cases.

Finally, convex duality gives the formula central to the new proof:
\begin{equation}\label{eq:dual}
 \psi(\xi,\rho)=
 \max_{0\le\lambda\le1}
 \left\{\lambda\xi-
 \frac{\lambda^2}{4(\rho+\lambda)}\right\}.
\end{equation}
Indeed, write $x_+=\max_{0\le\lambda\le1}\lambda x$, exchange min and max,
and minimise $(\rho+\lambda)v^2-\lambda v+\lambda\xi$ in $v$.

There are two canonical choices in \eqref{eq:dual}.  If $2\rho\xi\le1$, take
$\lambda=2\rho\xi$ to get
\begin{equation}\label{eq:certA}
 C_m\psi(\xi,\rho)
 \ge
 \frac{2\alpha^3A_m d^2}{e^2}
 \frac{1+4\xi}{1+2\xi}.
\end{equation}
If $2\rho\xi>1$, take $\lambda=1$.  Since
$1/[4(1+\rho)]<\xi/(2+4\xi)$ on this branch,
\begin{equation}\label{eq:certB}
 C_m\psi(\xi,\rho)
 \ge
 \sqrt{2\alpha}\,B_m f d
 \frac{1+4\xi}{2+4\xi}.
\end{equation}
These are the quadratic and linear certificates, respectively.

\section{The dimensionless frontier chart}\label{sec:chart}

Fix an admissible pair $(q,\alpha)$ and an odd $m\ge15$.  Put
\begin{equation}\label{eq:chart}
 N=m-2,\qquad
 a=\frac q\alpha,\qquad
 r=\frac Lq,\qquad
 \ell=\frac L\alpha=ar,
 \qquad z=\ell^2,\qquad u=1-a.
\end{equation}
Also set
\begin{equation}\label{eq:D-sigma}
 D=1+a^2+z,
 \qquad
 \sigma=\frac p\alpha=\frac{1+z}{a}.
\end{equation}
Dividing $L^2=pq-\alpha^2$ by $\alpha^2$ and using $p+q=1$ gives
\begin{equation}\label{eq:chart-identities}
 z=a\sigma-1,\qquad
 \alpha=\frac aD,\qquad q=\frac{a^2}{D},\qquad
 e=\frac{u^2+z}{D},\qquad
 \xi=\frac{4a^3u}{D(u^2+z)^2}.
\end{equation}

The structural conditions become especially simple:
\begin{equation}\label{eq:domain}
 0<r<1,\qquad
 z\ge au,\qquad
 D<3a^2,\qquad
 a>\frac{1+\sqrt{13}}6>\frac34.
\end{equation}
Here $z\ge au$ is exactly the frontier ceiling
$\alpha^2+q\alpha-q\le0$, and $D<3a^2$ is exactly $q>1/3$.
Their compatibility gives $3a^2-a-1>0$, proving the last bound.  We shall
also use
\begin{equation}\label{eq:domain-consequences}
        a^2(2-r^2)>1,\qquad a>r,\qquad
        \sigma>1+z.
\end{equation}

After division by $\alpha^m$, the scalar defect is
\begin{align}
 F_N:=\frac{R_m}{\alpha^m}
 &=1-(1+z)a^N+\ell^{N+2}\notag\\
 &=1-a^N-z(a^N-\ell^N).\label{eq:defect-normal}
\end{align}
Define the two finite quotients
\begin{equation}\label{eq:KA-KL}
 K_A=\frac{\sigma^{N+1}-1}{\sigma+1},
 \qquad
 K_L=\frac{\sigma^{N+1}-\ell^{N+1}}{\sigma+\ell}.
\end{equation}
Then
\begin{align}
 \frac{A_m}{\alpha^N}&=2\ell^N+(N+2)K_A,
 &
 \frac{B_m}{\alpha^N}&=2\ell^N+(N+2)K_L.
 \label{eq:AB-normal}
\end{align}

The next two sections prove that \eqref{eq:certA} and \eqref{eq:certB}, in
these coordinates, both dominate $F_N$ on their respective intrinsic
branches.

\section{The quadratic certificate}\label{sec:quadratic}

Assume \(2\rho\xi\le1\).  From \eqref{eq:certA} and
\(A_m\ge m k_m(\alpha)\),
\begin{equation}\label{eq:quadratic-start}
 \frac{C_m\psi(\xi,\rho)}{\alpha^m}
 \ge 2(N+2)\alpha\left(\frac d e\right)^2K_A.
\end{equation}
(The omitted factor \((1+4\xi)/(1+2\xi)\) is at least \(1\).)
The next lemma proves that the right side dominates the normalized defect.
It is valid on the whole admissible chart; the branch assumption is needed
only for the dual witness in \eqref{eq:quadratic-start}.

\begin{lemma}[Quadratic payment lemma]\label{lem:quadratic-payment}
For every point satisfying \eqref{eq:domain} and every odd \(N\ge13\),
\begin{equation}\label{eq:quadratic-lemma}
 F_N\le 2(N+2)\alpha\left(\frac d e\right)^2K_A.
\end{equation}
\end{lemma}

\begin{proof}
Introduce two further dimensionless variables
\begin{equation}\label{eq:zeta-v}
 \zeta=\frac{\ell^2}{u},\qquad v=\sigma-1.
\end{equation}
The frontier ceiling gives \(\zeta\ge a>3/4\), while
\(1<\sigma<2\) gives \(0<v<1\).  Solving the definitions yields
\begin{equation}\label{eq:zeta-param}
 H_\zeta:=\zeta+1+v,\qquad
 u=\frac v{H_\zeta},\qquad
 a=\frac{\zeta+1}{H_\zeta},\qquad
 \ell^2=\frac{\zeta v}{H_\zeta}.
\end{equation}
Set
\begin{equation}\label{eq:TN}
        T_N=N-\zeta(a^N-\ell^N).
\end{equation}
Since \(1-a^N\le Nu\) and \(\ell^2=\zeta u\),
\begin{equation}\label{eq:F-to-T}
        F_N\le uT_N.
\end{equation}
Thus we may assume \(T_N>0\).

Because \(N\) is odd, AM--GM in the finite geometric sum gives
\begin{equation}\label{eq:KA-amgm}
 K_A=(\sigma-1)\sum_{j=0}^{(N-1)/2}\sigma^{2j}
 \ge\frac{N+1}{2}(\sigma-1)\sigma^{(N-1)/2}.
\end{equation}
Also
\[
 \alpha=\frac aD,\qquad
 \frac d e=\frac{a}{u+\zeta},\qquad
 \sigma-1=\frac{u(1+\zeta)}a.
\]
Consequently the right side of \eqref{eq:quadratic-lemma} is at least
\begin{equation}\label{eq:quadratic-payment-lower}
 u\,\frac{(N+2)(N+1)a^2(1+\zeta)\sigma^{(N-1)/2}}
          {D(u+\zeta)^2}.
\end{equation}

We need one coefficient estimate:
\begin{equation}\label{eq:coefficient-lemma}
 \frac{a^2\sigma^{5/2}}{D(u+\zeta)^2}\ge\frac1{2\zeta^2}.
\end{equation}
Indeed, \(u+\zeta\le\zeta/a\), \(D=a(\sigma+a)\), and it is enough to
show \(2a^3\sigma^{5/2}\ge\sigma+a\).  The left side minus the right side
is increasing in \(\sigma\).  Moreover,
\[
 \sigma\ge\sigma_0:=\frac{1+a-a^2}{a}.
\]
Since \(\sqrt{\sigma_0}\ge1+(\sigma_0-1)/3\), substitution reduces the
desired inequality to
\[
 \frac{(1-a)P(a)}{3a}>0,\qquad
 P(a)=2a^6-8a^5+14a^3+6a^2-4a-3.
\]
Here
\[
 P'(a)=2Q(a),\qquad
 Q(a)=6a^5-20a^4+21a^2+6a-2
 \ge a^2+6a-2>0
\]
on \(a\ge3/4\), and \(P(3/4)=3561/2048>0\).  This proves
\eqref{eq:coefficient-lemma}.

Put
\begin{equation}\label{eq:Czeta}
 C_N(\zeta)=\frac{(N+2)(N+1)(1+\zeta)}{2\zeta^2},
 \qquad R=\frac{N-6}{2}.
\end{equation}
Equations \eqref{eq:quadratic-payment-lower} and
\eqref{eq:coefficient-lemma} show that it remains only to prove
\begin{equation}\label{eq:TN-core}
        T_N\le C_N(\zeta)(1+v)^R.
\end{equation}
This estimate has a natural interpretation: \(\zeta\) measures frontier
variance in units of the gap \(u\), and the comparison changes when that
variance crosses the cycle scale \(N\).

From \eqref{eq:zeta-param}, Bernoulli's inequality gives the two bounds
\begin{align}
 T_N&\le N-\zeta+Nv+\zeta v^{N/2},\label{eq:T-first}\\
 T_N&<N(1-\ell^{N+1}).\label{eq:T-second}
\end{align}
For the second, note that
\(\ell^2<2a^2-1<(2a-1)^2\), so \(a-\ell>1-a=u\), and use
convexity: \(a^N-\ell^N>N u\ell^{N-1}\).

\smallskip
\noindent\emph{Variance below the cycle scale.}
If \(\zeta\le N/2\), then \(C_N(\zeta)\ge C_N(N/2)>N>T_N\).
If \(N/2\le\zeta\le N\), then
\begin{equation}\label{eq:C-two}
        C_N(\zeta)\ge N-\zeta,\qquad
        C_N(\zeta)R\ge N+\zeta.
\end{equation}
The first inequality follows from
\(C_N(\zeta)\ge C_N(N)>N/2\); for the second, the worst endpoint is
\(\zeta=N\), where it is equivalent to
\[
 (N+2)(N+1)^2(N-6)\ge8N^3,
\]
which holds for \(N\ge13\).  Bernoulli's inequality and
\(v^{N/2}\le v\) now give
\[
 C_N(\zeta)(1+v)^R
 \ge(N-\zeta)+(N+\zeta)v
 \ge T_N.
\]

\smallskip
\noindent\emph{Variance above the cycle scale.}
Suppose \(\zeta\ge N\).  Positivity of \(T_N\), together with
\(a^N\ge1-Nv/\zeta\) and \(\ell^N\le v^{N/2}\), implies
\begin{equation}\label{eq:zeta-upper}
 \zeta<Z:=\frac{N(1+v)}{1-v^{N/2}}.
\end{equation}

First let \(0<v\le1/2\).  From \eqref{eq:T-first},
\[
        T_N\le N(v+v^{N/2}).
\]
Using \eqref{eq:zeta-upper} and \(1+\zeta>\zeta\),
\begin{align}
 C_N(\zeta)(1+v)^R
 &>
 \frac{(N+2)(N+1)}{2N}
 (1-v^{N/2})(1+v)^{R-1}.\label{eq:Mv}
\end{align}
Now \(v^{N/2}\le v/32\), \(1-v^{N/2}>63/64\), and
\((1+v)^{R-1}\ge1+5v/2\).  The right side of \eqref{eq:Mv} is therefore
larger than \(33Nv/32\), because
\[
 \frac{63(N+1)}{128}\left(1+\frac52v\right)
 -\frac{33N}{32}v>0.
\]
This dominates \(N(v+v^{N/2})\).

It remains to consider \(1/2\le v<1\).  We use the vanishing in
\eqref{eq:T-second} more accurately.  For fixed \(v\), define
\[
 g(t)=t\left[1-
 \left(\frac{tv}{t+1+v}\right)^{(N+1)/2}\right].
\]
Writing \(\theta=(1+v)/(t+1+v)\), differentiation gives
\[
 g'(t)=1-[v(1-\theta)]^{(N+1)/2}
 \left(1+\frac{N+1}{2}\theta\right)\ge0;
\]
the last inequality is Bernoulli's inequality in the form
\((1-\theta)^{-(N+1)/2}\ge1+(N+1)\theta/2\).
Hence \eqref{eq:T-second} and \eqref{eq:zeta-upper} imply
\begin{equation}\label{eq:gZ}
 \zeta T_N<N g(\zeta)<N g(Z)
 =\frac{N^2(1+v)}{1-y}(1-\omega^{(N+1)/2}),
\end{equation}
where
\[
        y=v^{N/2},\qquad
        \omega=\frac{Nv}{N+1-y}.
\]
The following tangent estimate is uniform:
\begin{equation}\label{eq:tangent}
 1-\omega^{(N+1)/2}
 \le\frac{3(N+1)}{2N}(1-y).
\end{equation}
Indeed, as a function of \(y\),
\[
 f(y):=\omega^{(N+1)/2}
 =\left(\frac{N}{N+1-y}\right)^{(N+1)/2}
  y^{(N+1)/N}
\]
satisfies \(f(1)=1\), \(f'(1)=3(N+1)/(2N)\), and \(f''(y)>0\).
Convexity places \(f\) above its tangent at \(1\), proving
\eqref{eq:tangent}.  Equations \eqref{eq:gZ} and \eqref{eq:tangent} give
\[
        T_N<\frac{3N(N+1)}{2\zeta}(1+v).
\]
On the other hand,
\[
 C_N(\zeta)(1+v)^R>
 \frac{(N+2)(N+1)}{2\zeta}(1+v)^R.
\]
Thus \eqref{eq:TN-core} follows from
\[
 (1+v)^{R-1}\ge\frac{3N}{N+2}.
\]
For \(N=13\), this is implied by
\((3/2)^{5/2}>13/5\), whose square is
\(243/32>169/25\).  For odd \(N\ge15\), the left side is at least
\((3/2)^{7/2}>3\), while the right side is smaller than \(3\).
This completes the proof of \eqref{eq:TN-core} and of the lemma.
\end{proof}

\begin{proof}[Proof of Proposition~\ref{prop:quadratic}]
Combine \eqref{eq:quadratic-start} with
Lemma~\ref{lem:quadratic-payment} and \(F_N=R_m/\alpha^m\).
\end{proof}

\section{The linear certificate}\label{sec:linear}

Assume now that $2\rho\xi>1$, and abbreviate
\begin{equation}\label{eq:varphi-c}
        \varphi=\frac{1+4\xi}{2+4\xi},
        \qquad c_\xi=\sqrt{2\alpha}\,\varphi.
\end{equation}
By \eqref{eq:certB} and \eqref{eq:AB-normal},
\begin{equation}\label{eq:linear-start}
 \frac{C_m\psi(\xi,\rho)}{\alpha^m}
 \ge u(N+2)c_\xi(1-\ell)K_L.
\end{equation}
Since $N+1$ is even,
\begin{equation}\label{eq:KL-lower}
 K_L=(\sigma-\ell)
 \sum_{j=0}^{(N-1)/2}\sigma^{N-1-2j}\ell^{2j}
 \ge(\sigma-\ell)\sigma^{N-1}.
\end{equation}
Recall the variables $\zeta,v,T_N$ from
\eqref{eq:zeta-v}--\eqref{eq:TN}.  In view of
$F_N\le uT_N$, Proposition~\ref{prop:linear} will follow from
\begin{equation}\label{eq:linear-core}
 T_N\le
 (N+2)c_\xi(1-\ell)(\sigma-\ell)\sigma^{N-1}.
\end{equation}

The key point is that the Huber factor $\varphi$ compensates automatically
for the possible loss when $\ell$ is close to $1$.

\begin{lemma}[Compensation]\label{lem:compensation}
On the whole admissible chart,
\begin{equation}\label{eq:compensation}
        c_\xi\,\sigma^{\zeta/4}\ge1-\ell^2.
\end{equation}
\end{lemma}

\begin{proof}
We prove separately the two factors in $c_\xi$.
First,
\begin{equation}\label{eq:sqrt-comp}
        \sqrt{2\alpha}\ge1-\ell^2.
\end{equation}
Indeed, with $y=\ell^2$, this is equivalent after squaring to
\[
        2a\ge(1+a^2+y)(1-y)^2.
\]
The difference between the two sides is increasing in $y$, since its
derivative is $(1-y)(2a^2+3y+1)$.  At the smallest allowed value
$y=au$ it equals
\[
        (1-a)(a^4+a^2+2a-1)>0.
\]

It remains to prove $\varphi\sigma^{\zeta/4}\ge1$.  In terms of
$H_\zeta=\zeta+1+v$, direct substitution gives
\begin{equation}\label{eq:xi-zeta-v}
 \xi=\frac{4H_\zeta^2}
 {v(\zeta+v)^2Q},
 \qquad
 Q=v^2+v\zeta+2v+2\zeta+2<3(\zeta+1).
\end{equation}
Moreover,
\[
 \zeta H_\zeta^2-(\zeta+1)(\zeta+v)^2
 =\zeta^2+\zeta-v^2>0.
\]
Consequently
\begin{equation}\label{eq:xi-large-enough}
        4\zeta v\xi>\frac{16}{3}.
\end{equation}
Since $\zeta>3/4$ and $0<v<1$, this implies
\[
        \zeta v(1+4\xi)>2(v+2).
\]
(For $\zeta\ge2$ this is immediate; for $\zeta<2$, use
$v(2-\zeta)<5/4<4/3$.)  Finally,
\[
 \log\sigma=\log(1+v)\ge\frac{2v}{v+2},
\]
and hence
\[
 \frac\zeta4\log\sigma>
 \frac1{1+4\xi}.
\]
Exponentiating and using $e^t>1+t$ gives
\[
 \sigma^{\zeta/4}>
 1+\frac1{1+4\xi}=\frac1\varphi.
\]
Together with \eqref{eq:sqrt-comp}, this proves
\eqref{eq:compensation}.
\end{proof}

We shall also need one elementary one-variable inequality.

\begin{lemma}[Geometric growth]\label{lem:geometric-growth}
Let $N\ge13$ and put $A=3N/4-1$.  If $0\le x<1$ and
\begin{equation}\label{eq:growth-assumption}
        1-x\ge\frac{2x^2}{N},
\end{equation}
then
\begin{equation}\label{eq:growth-conclusion}
        (1-x)^2(1+x^2)^A\ge\frac{N}{N+2}.
\end{equation}
\end{lemma}

\begin{proof}
Let $G(x)=(1-x)^2(1+x^2)^A$, and let $\beta$ be the positive solution of
$1-\beta=2\beta^2/N$.  Condition \eqref{eq:growth-assumption} says
$0\le x\le\beta$.  The sign of the logarithmic derivative of $G$ is the
sign of
\begin{equation}\label{eq:growth-derivative}
        -(A+1)x^2+Ax-1.
\end{equation}
Thus a global minimum occurs at $0$, at $\beta$, or at the smaller critical
point $t$ of \eqref{eq:growth-derivative}, if it exists.

At $x=0$, $G=1$.  At a critical point,
$At(1-t)=1+t^2$, and Bernoulli's inequality gives
\[
        G(t)\ge(1-t)^2(1+At^2)=(1-t)(1+t^3).
\]
The smaller critical point satisfies $0<t\le7/52$.  Eliminating $N$ via the
critical-point equation gives
\begin{equation}\label{eq:growth-critical}
 (1-t)(1+t^3)-\frac{N}{N+2}
 =\frac{t(1+t)P(t)}{(2-t)(1+3t)},
\end{equation}
where
\[
        P(t)=3t^4-11t^3+14t^2-9t+1.
\]
The polynomial $P$ is decreasing on $[0,7/52]$ and
$P(7/52)=119263/7311616>0$, so \eqref{eq:growth-critical} is nonnegative.

Finally, $\beta\ge7/8$, and therefore
\begin{align*}
 G(\beta)
 &=\frac{4\beta^4}{N^2}(1+\beta^2)^A\\
 &\ge
 \frac{4(7/8)^4}{N^2}
 \left(\frac{113}{64}\right)^{3N/4-1}>1.
\end{align*}
The last inequality holds at $N=13$ even with the exponent replaced by $8$;
its ratio under $N\mapsto N+2$ is larger than $1$.  Thus every possible
minimum is at least $N/(N+2)$.
\end{proof}

\begin{lemma}[Linear payment lemma]\label{lem:linear-payment}
For every point satisfying \eqref{eq:domain} and every odd $N\ge13$,
inequality \eqref{eq:linear-core} holds.
\end{lemma}

\begin{proof}
If $T_N\le0$ the claim is immediate, so assume $T_N>0$.

\smallskip
\noindent\emph{Variance below the cycle scale: $\zeta\le N$.}
By Lemma~\ref{lem:compensation},
\begin{align}
 &(N+2)c_\xi(1-\ell)(\sigma-\ell)\sigma^{N-1}\notag\\
 &\quad\ge
 (N+2)(1-\ell)^2(1+\ell^3)
 \sigma^{N-1-\zeta/4}\notag\\
 &\quad\ge
 (N+2)(1-\ell)^2(1+\ell^2)^{3N/4-1}.
 \label{eq:linear-below}
\end{align}
Here we used $\sigma\ge1+\ell^2$ and
$(1+\ell)(1-\ell+\ell^2)=1+\ell^3$.
The domain inequality $\ell^2<2a^2-1$ implies
$\ell<2a-1$, whence
\[
 1-\ell>2u=\frac{2\ell^2}{\zeta}
 \ge\frac{2\ell^2}{N}.
\]
Lemma~\ref{lem:geometric-growth}, with $x=\ell$, shows that the final
line of \eqref{eq:linear-below} is at least $N$.  Since
$T_N=N-\zeta(a^N-\ell^N)<N$, this proves the claim.

\smallskip
\noindent\emph{Variance above the cycle scale: $\zeta\ge N$.}
Recall $v=\sigma-1$.  From \eqref{eq:zeta-param},
$\ell^2<v$, and \eqref{eq:T-first} gives
\begin{equation}\label{eq:T-high-zeta}
        T_N\le N(v+v^{N/2}).
\end{equation}
Also $c_\xi>1/\sqrt6$, because $\alpha>1/3$ and $\varphi>1/2$.

Suppose first that $0<v\le1/2$, and put $y=\sqrt v$.  Since
$\ell\le y$,
\begin{equation}\label{eq:Py}
        (1-\ell)(\sigma-\ell)
        \ge(1-y)(1-y+y^2)=:P_y.
\end{equation}
If $v\le1/4$, then $P_y\ge3/8$,
$v+v^{N/2}\le2v$, and
$(1+v)^{N-1}\ge1+(N-1)v$.  Thus \eqref{eq:linear-core} follows from
\[
        \frac{2N}{N-1}\le\frac{3(N+2)}{8\sqrt6},
\]
which holds at $N=13$ and becomes stronger as $N$ increases.
If $1/4\le v\le1/2$, then $P_y\ge3/16$ and the left side of
\eqref{eq:T-high-zeta} is at most $N$.  It is enough that
\[
        \frac{3(N+2)}{16\sqrt6}
        \left(\frac54\right)^{N-1}>N.
\]
Again the inequality holds at $N=13$ (squaring clears the only radical) and
its ratio increases with $N$.

Finally suppose $v\ge1/2$.  From \eqref{eq:T-second},
\[
        T_N<N(N+1)(1-\ell).
\]
On the other hand, $\sigma-\ell>v\ge1/2$ and $\sigma\ge3/2$.
After cancelling $1-\ell$, it is enough to prove
\begin{equation}\label{eq:last-growth}
        \left(\frac32\right)^{N-1}
        \ge\frac{2\sqrt6\,N(N+1)}{N+2}.
\end{equation}
For $N\ge13$, one has $(3/2)^{N-1}\ge5N$: check $N=13$ and use the
ratio under $N\mapsto N+1$.  Since
$5>2\sqrt6\,(N+1)/(N+2)$, this proves
\eqref{eq:last-growth}.  All values of $\zeta$ and $v$ have now been
covered.
\end{proof}

\begin{proof}[Proof of Proposition~\ref{prop:linear}]
Combine \eqref{eq:linear-start}, \eqref{eq:KL-lower}, and
Lemma~\ref{lem:linear-payment}, then use $F_N=R_m/\alpha^m$.
\end{proof}

\section{Assembly and scope}\label{sec:assembly}

\begin{proof}[Proof of Theorem~\ref{thm:main}]
Let $q=1-p\in(1/3,1/2)$.  If the compression $A$ has no eigenvalue larger
than $q$, the no-frontier argument following \eqref{eq:shift} proves
\eqref{eq:goal}.  Otherwise there is a unique frontier eigenvalue
$q<\alpha\le r(q)$, and \eqref{eq:master}--\eqref{eq:huber-elim} reduce the
claim to
\[
        R_m\le C_m\psi(\xi,\rho).
\]
If $2\rho\xi\le1$, this is Proposition~\ref{prop:quadratic}; if
$2\rho\xi>1$, it is Proposition~\ref{prop:linear}.  These alternatives are
exhaustive.
\end{proof}

\begin{remark}[Nature of the remaining subdivisions]
The proof of each scalar lemma compares a geometric parameter with its
natural scale: $\zeta$ with the cycle length $N$, and $v=\sigma-1$ with a
fixed distance from the endpoint $1$.  These subdivisions only select among
elementary estimates of the same scalar expression.  They do not create
different graphon configurations, different spectral reductions, or
separately certified parameter boxes.
\end{remark}

\begin{remark}[Short odd cycles]
The threshold $N\ge13$, equivalently $m\ge15$, enters only through coarse
uniform estimates such as Lemma~\ref{lem:geometric-growth}.  Direct numerical
inspection of the exact scalar inequality does not suggest a failure for
$m\le13$; rather, these particular lower bounds cease to close.  A genuinely
uniform computer-free proof of $C_5,\ldots,C_{13}$ would require a sharper
finite-cycle lemma and is not supplied here.  Thus this document should be
combined with the existing per-cycle results only when stating the full
conjecture.
\end{remark}

\begin{remark}[Auditability]
Every non-symbolic endpoint check in this proof is displayed as a rational
inequality or can be reduced to one by squaring a positive radical.  No claim
depends on floating-point evaluation.  The proof is therefore suitable for a
conventional line-by-line audit independently of the former certifier code.
\end{remark}

\end{document}
